\section{Running the Application}
\label{sec:execution}

\subsection{The Two-Terminal Execution Model}
Research Rover operates as an asynchronous decoupled architecture. The backend API server and the frontend client interface must run concurrently in two separate terminal sessions.

\begin{figure}[H]
\centering
\begin{minipage}{0.48\textwidth}
\textbf{Terminal 1: FastAPI Backend}
\begin{lstlisting}[language=bash]
cd E:\Research_Rover\backend
.venv\Scripts\Activate.ps1
uvicorn main:app --host 0.0.0.0 --port 8000 --reload
\end{lstlisting}
\end{minipage}
\hfill
\begin{minipage}{0.48\textwidth}
\textbf{Terminal 2: Vite Frontend}
\begin{lstlisting}[language=bash]
cd E:\Research_Rover\frontend
npm run dev
\end{lstlisting}
\end{minipage}
\caption{Two-Terminal Execution Commands for Running Research Rover Locally}
\label{fig:two_terminal_box}
\end{figure}

\subsection{Starting the Services}

\subsubsection{Terminal 1: Start the Backend Server}
In your first terminal, activate the Python virtual environment and launch Uvicorn:

\begin{lstlisting}[language=bash, caption={Launching FastAPI Backend Server}]
cd E:\Research_Rover\backend
.venv\Scripts\Activate.ps1
uvicorn main:app --host 0.0.0.0 --port 8000 --reload
\end{lstlisting}

\noindent Upon successful startup, the terminal logs will display:
\begin{lstlisting}
INFO:     Started server process [pid]
INFO:     Waiting for application startup.
INFO:     Starting Research Rover API...
INFO:     Research Rover API started successfully
INFO:     Uvicorn running on http://0.0.0.0:8000 (Press CTRL+C to quit)
\end{lstlisting}

\noindent The application lifespan context manager initializes the data directory at \texttt{E:\textbackslash Research\_Rover\textbackslash data} and begins asynchronous background warm-up of transformer models without blocking incoming HTTP traffic.

\subsubsection{Terminal 2: Start the Frontend Development Server}
In your second terminal, navigate to the frontend directory and start Vite:

\begin{lstlisting}[language=bash, caption={Launching Vite Frontend Server}]
cd E:\Research_Rover\frontend
npm run dev
\end{lstlisting}

\noindent The console output will confirm:
\begin{lstlisting}
  VITE v6.2.0  ready in 480 ms

  ➜  Local:   http://localhost:3000/
  ➜  Network: use --host to expose
  ➜  press h + enter to show help
\end{lstlisting}

\subsection{Application Access URLs and Health Verification}
Once both servers are operational, access the application and documentation interfaces via any standard web browser using the URLs listed in Table~\ref{tab:access_urls}.

\begin{table}[H]
\centering
\caption{Research Rover Local Access URLs and Service Interfaces}
\label{tab:access_urls}
\begin{tabularx}{\textwidth}{l X l}
\toprule
\textbf{Interface} & \textbf{Network URL} & \textbf{Primary Purpose} \\
\midrule
Web Application UI & \url{http://localhost:3000} & Interactive research portal and dashboards \\
Backend API Base & \url{http://localhost:8000} & Core REST API service \\
OpenAPI Swagger UI & \url{http://localhost:8000/docs} & Interactive API testing and parameter explorer \\
ReDoc Documentation & \url{http://localhost:8000/redoc} & Formal REST API specification documentation \\
System Health Check & \url{http://localhost:8000/health} & Real-time model and runtime status telemetry \\
Root Ping Endpoint & \url{http://localhost:8000/} & Lightweight health check (\texttt{\{"status": "healthy"\}}) \\
\bottomrule
\end{tabularx}
\end{table}

\subsection{Terminating the Application}
To shut down the platform cleanly:
\begin{enumerate}
    \item Focus on Terminal 2 (Frontend) and press \texttt{Ctrl + C}. Confirm termination if prompted.
    \item Focus on Terminal 1 (Backend) and press \texttt{Ctrl + C}. Uvicorn will execute the lifespan shutdown sequence and release bound network ports.
\end{enumerate}
